\origpage{352}{N. Robertson, P.D. Seymour / Journal of Combinatorial Theory, Series B 90 (2004) 325--385}
If $\chi$ is disjoint, $(\Gamma, \phi)$ is a $\chi$-coloured painting, and $\mathscr{L}$ is a rooted location in $\Gamma$, we call $(\Gamma, \phi, \mathscr{L})$ a $\chi$-\textit{place}. A set $\mathscr{P}$ of $\chi$-places is \textit{well-behaved} if for every well-quasi-order $\Omega$ and every countable sequence $(\Gamma_i, \phi_i, \mathscr{L}_i)$ ($i = 1, 2, \ldots$) of members of $\mathscr{P}$, and for all functions $\xi_i \colon \mathscr{L}_i \to E(\Omega)$, there exist $j > i \ge 1$ and an outline $\tau$ of $(\Gamma_i, \phi_i, \mathscr{L}_i)$ in $(\Gamma_j, \phi_j, \mathscr{L}_j)$ such that $\xi_i(A) \le \xi_j(\tau(A))$ for all $A \in \mathscr{L}_i$. Theorem 6.7 of [10], together with 9.1, imply the following.

\textbf{9.2.} \textit{Let $\chi$ be a disjoint colour scheme, let $(\Gamma_i, \phi_i)$ ($i = 1, 2, \ldots$) be a countable sequence of $\chi$-coloured paintings and for each $i \ge 1$ let $\lambda_i$ be a tie-breaker in $\Gamma_i$. Let $\theta \ge 1$, and for each $i \ge 1$ and every tangle $\mathscr{T}$ in $\Gamma_i$ of order $\ge \theta$, let $\mathscr{L}(\mathscr{T})$ be a rooted location in $\Gamma_i$ which $\theta$-isolates $\mathscr{T}$. Let the set of all these $(\Gamma_i, \phi_i, \mathscr{L}(\mathscr{T}))$ (over all $i$ and $\mathscr{T}$) be well-behaved. Then there exist $j > i \ge 1$ such that there is an inflation of $(\Gamma_i, \phi_i)$ in $(\Gamma_j, \phi_j)$.}

Now we can give the reader a little better intuition as to how the proof of 4.1 will work. Let $\chi$ be a colour scheme satisfying $\mathbf{S}_1, \ldots, \mathbf{S}_4$, and let $(\Gamma_i, \phi_i)$ ($i = 1, 2, \ldots$) be a countable sequence of $\chi$-coloured paintings. By 6.1, it will suffice to show that there exist $j > i \ge 1$ such that there is an inflation of $(\Gamma_i, \phi_i)$ in $(\Gamma_j, \phi_j)$, and to show this we will apply 9.2. We therefore need to produce an infinite subsequence of this sequence, and a well-behaved set of rooted locations, so that the subsequence satisfies the hypotheses of 9.2. To get the subsequence, discard from the given sequence the first term, and all terms $(\Gamma_i, \phi_i)$ with $dist(\Gamma_i)$ or $rep(\Gamma_i)$ at most some appropriately-chosen number that depends only on $(\Gamma_1, \phi_1)$. By 5.4 and 5.5, an infinite sequence still remains, and this is the one we need. Now we need to produce the well-behaved set of rooted locations. Let $(\Gamma_i, \phi_i)$ be some term of the sequence that still remains. We know that there is no inflation of $(\Gamma_1, \phi_1)$ in $(\Gamma_i, \phi_i)$; and a theorem of [7] therefore can be applied. That theorem implies that for every large-order tangle $\mathscr{T}$ in $\Gamma_i$ (the meaning of ``large'' depending only on $(\Gamma_1, \phi_1)$), the triple $(\Gamma_i, \phi_i, \mathscr{T})$ is ``insufficiently general'', in one of only a few possible ways. We deduce that there is a ``flaw'' in $(\Gamma_i, \phi_i, \mathscr{T})$, of one of only a few possible kinds (and in particular, of only finitely many different kinds). For each kind of flaw, we shall show that there corresponds a well-behaved set of rooted locations, such that if $(\Gamma_i, \phi_i, \mathscr{T})$ admits the flaw then some rooted location in this set $\theta$-isolates $\mathscr{T}$ (for appropriate $\theta$); and then the union of these finitely many well-behaved sets is another well-behaved set, now satisfying the hypotheses of 9.2, as required.

In Sections 10--14 we look at the various kinds of flaw, and in each case construct the desired well-behaved set, and in Section 15 we complete the proof by applying the theorem of [7].

We shall need the following lemmas.

\textbf{9.3.} \textit{Let $\chi$ be a disjoint colour scheme, and let $\mathscr{P}$ be a set of $\chi$-places. For each $(\Gamma, \phi, \mathscr{L}) \in \mathscr{P}$ let $\pi(\mathscr{L})$ be a march with $\bar{\pi}(\mathscr{L}) = \bigcup (\bar{\pi}(A) : A \in \mathscr{L})$. Suppose that}

\hspace*{12pt}$\bullet$ \textit{there exist $m, n$ such that $|\mathscr{L}| \le m$ and $|\bar{\pi}(\mathscr{L})| \le n$ for all $(\Gamma, \phi, \mathscr{L}) \in \mathscr{P}$,}
